Un lait frais ne contient pas d'acide lactique. En vieillissant, le lactose
présent dans le lait se transforme en acide lactique, noté par la suite
$\ce{HA}$.

On fait réagir un volume $V_{A}=\qty{20}{\mL}$ de lait contenant de l'acide
lactique, considéré comme le seul acide présent dans le lait étudié, avec une
solution d'hydroxyde de sodium $\ce{Na^+_{(aq)} + HO^-_{(aq)}}$ (soude) de
concentration $C_{B}=\qty{5,00e-2}{\M}$.

\begin{donnees}
    \begin{itemize}
        \item pour le couple $\ce{HA_{(aq)}/A^-_{(aq)}}$~:
              $\mathrm{pK_{A}}=3,9$.
    \end{itemize}
\end{donnees}

\begin{enumerate}
  \item Écrire l'équation de la réaction qui se produit lors du mélange.
  \item Exprimer la constante de réaction $\mathrm{K}$ correspondante.
  \item Le $\mathrm{pH}$ du lait est de $2,9$. En utilisant un diagramme de
        prédominance, déterminer quelle est, entre $\ce{HA_{(aq)}}$ et
        $\ce{A^-_{(aq)}}$, l'espèce chimique prédominante.
\end{enumerate}

